\section{Тестирование программной системы}

\subsection{Цели и задачи тестирования}

Целью тестирования является проверка соответствия разработанного приложения требованиям, сформулированным в техническом задании, а также выявление и устранение дефектов.

Основные задачи тестирования:
\begin{enumerate}
	\item проверка корректности реализации функциональных требований;
	\item проверка устойчивости приложения к некорректному вводу данных;
	\item проверка производительности приложения;
	\item проверка работы приложения в условиях ограниченных ресурсов;
	\item проверка совместимости с различными версиями Android.
\end{enumerate}

\subsection{План тестирования}

Тестирование проводилось на следующих устройствах:
\begin{itemize}
	\item эмулятор Android Pixel 4 (Android 11, API 30);
	\item эмулятор Android Pixel 6 (Android 13, API 33);
	\item физическое устройство Samsung Galaxy A51 (Android 12);
	\item физическое устройство Xiaomi Redmi Note 9 (Android 11);
	\item физическое устройство Google Pixel 3 (Android 10).
\end{itemize}

\subsection{Функциональное тестирование}

В таблице \ref{tab:test_cases} приведены тестовые сценарии для проверки функциональных требований.

\renewcommand{\arraystretch}{1.3}
\begin{longtable}{|p{1.4cm}|p{4cm}|p{3.8cm}|p{3cm}|p{2.2cm}|}
	\caption{Функциональные тестовые сценарии}
	\label{tab:test_cases} \\
	\hline
	\textbf{ID} & \textbf{Сценарий} & \textbf{Ожидаемый результат} & \textbf{Фактический результат} & \textbf{Статус} \\
	\hline
	\endfirsthead
	
	\caption[]{Продолжение таблицы \ref{tab:test_cases}} \\
	\hline
	\textbf{ID} & \textbf{Сценарий} & \textbf{Ожидаемый результат} & \textbf{Фактический результат} & \textbf{Статус} \\
	\hline
	\endhead
	
	\hline
	\multicolumn{5}{r}{\textit{Продолжение на следующей странице}} \\
	\endfoot
	
	\hline
	\endlastfoot
	
	TC-01 & Регистрация нового пользователя с корректными данными & Пользователь зарегистрирован, выполнен автоматический вход & Соответствует & Пройдено \\
	\hline
	TC-02 & Регистрация с уже существующим email & Сообщение «Пользователь уже существует» & Соответствует & Пройдено \\
	\hline
	TC-03 & Регистрация с пустыми полями & Подсветка полей, сообщение об ошибке & Соответствует & Пройдено \\
	\hline
	TC-04 & Авторизация с корректными данными & Вход выполнен, открыт главный экран & Соответствует & Пройдено \\
	\hline
	TC-05 & Авторизация с неверным паролем & Сообщение «Неверный email или пароль» & Соответствует & Пройдено \\
	\hline
	TC-06 & Авторизация с несуществующим email & Сообщение «Неверный email или пароль» & Соответствует & Пройдено \\
	\hline
	TC-07 & Просмотр каталога занятий & Список занятий отображается & Соответствует & Пройдено \\
	\hline
	TC-08 & Фильтрация занятий по городу & Отображаются только занятия выбранного города & Соответствует & Пройдено \\
	\hline
	TC-09 & Просмотр детальной информации о занятии & Открывается экран с полной информацией & Соответствует & Пройдено \\
	\hline
	TC-10 & Запись на занятие (есть места) & Запись сохранена, уведомление получено & Соответствует & Пройдено \\
	\hline
	TC-11 & Запись на занятие (нет мест) & Сообщение «Нет свободных мест» & Соответствует & Пройдено \\
	\hline
	TC-12 & Запись на одно занятие дважды & Сообщение «Вы уже записаны» & Соответствует & Пройдено \\
	\hline
	TC-13 & Отмена записи & Запись удалена, место освобождено & Соответствует & Пройдено \\
	\hline
	TC-14 & Просмотр расписания с календарём & Отображаются занятия на выбранную дату & Соответствует & Пройдено \\
	\hline
	TC-15 & Просмотр списка инструкторов & Список инструкторов отображается & Соответствует & Пройдено \\
	\hline
	TC-16 & Редактирование профиля & Данные обновлены в SharedPreferences & Соответствует & Пройдено \\
	\hline
	TC-17 & Загрузка аватара & Изображение отображается в профиле & Соответствует & Пройдено \\
	\hline
	TC-18 & Просмотр прогресса в виде графиков & Графики отображаются корректно & Соответствует & Пройдено \\
	\hline
	TC-19 & Оставление отзыва (рейтинг 5 звёзд) & Отзыв сохранён, рейтинг обновлён & Соответствует & Пройдено \\
	\hline
	TC-20 & Оставление отзыва без текста & Отзыв сохранён (текст не обязателен) & Соответствует & Пройдено \\
	\hline
	TC-21 & Просмотр списка активных записей & Список записей отображается & Соответствует & Пройдено \\
	\hline
	TC-22 & Получение напоминания о занятии & Уведомление получено за 1 час & Соответствует & Пройдено \\
	\hline
\end{longtable}
\renewcommand{\arraystretch}{1}

\subsection{Тестирование пользовательского интерфейса}

В таблице \ref{tab:ui_tests} приведены результаты тестирования пользовательского интерфейса.

\renewcommand{\arraystretch}{1.3}
\begin{table}[H]
	\centering
	\caption{Результаты тестирования пользовательского интерфейса}
	\label{tab:ui_tests}
	\normalsize
	\setlength{\tabcolsep}{6pt}
	\begin{tabular}{|p{5cm}|p{4.2cm}|p{3.3cm}|p{2.5cm}|}
		\hline
		\textbf{Проверяемый элемент} & \textbf{Ожидаемый результат} & \textbf{Фактический результат} & \textbf{Статус} \\
		\hline
		BottomNavigationView & Переключение между 5 вкладками & Соответствует & Пройдено \\
		\hline
		Календарь & Выбор даты, подсветка выбранного дня & Соответствует & Пройдено \\
		\hline
		RecyclerView (список занятий) & Плавная прокрутка, загрузка изображений & Соответствует & Пройдено \\
		\hline
		RatingBar (рейтинг) & Выбор от 1 до 5 звёзд & Соответствует & Пройдено \\
		\hline
		Spinner (выбор города) & Раскрывающийся список с 3 городами & Соответствует & Пройдено \\
		\hline
		Диалог редактирования профиля & Открытие формы с текущими данными & Соответствует & Пройдено \\
		\hline
		Модальное окно с фото & Увеличенное изображение, затемнение фона & Соответствует & Пройдено \\
		\hline
	\end{tabular}
\end{table}
\renewcommand{\arraystretch}{1}

\subsection{Тестирование производительности}

В таблице \ref{tab:performance} приведены результаты тестирования производительности.

\renewcommand{\arraystretch}{1.3}
\begin{table}[H]
	\centering
	\caption{Результаты тестирования производительности}
	\label{tab:performance}
	\normalsize
	\setlength{\tabcolsep}{6pt}
	\begin{tabular}{|p{5cm}|p{3cm}|p{3cm}|p{2.5cm}|}
		\hline
		\textbf{Сценарий} & \textbf{Требование} & \textbf{Фактическое значение} & \textbf{Статус} \\
		\hline
		Время запуска приложения & < 3 секунд & 1,8 сек & Пройдено \\
		\hline
		Время загрузки списка занятий & < 1 секунды & 0,3 сек & Пройдено \\
		\hline
		Время отклика на нажатие кнопки & < 0,5 секунды & 0,1 сек & Пройдено \\
		\hline
		Потребление ОЗУ (в фоне) & < 100 МБ & 42 МБ & Пройдено \\
		\hline
		Потребление ОЗУ (активное использование) & < 200 МБ & 78 МБ & Пройдено \\
		\hline
		Размер APK-файла & < 50 МБ & 8,2 МБ & Пройдено \\
		\hline
		Плавность прокрутки RecyclerView & 60 FPS & 60 FPS & Пройдено \\
		\hline
	\end{tabular}
\end{table}
\renewcommand{\arraystretch}{1}

\subsection{Тестирование в условиях ограниченных ресурсов}

Приложение было протестировано в следующих условиях:
\begin{itemize}
	\item отсутствие свободного места на устройстве — приложение корректно отображает сообщение об ошибке;
	\item прерывание работы приложения (входящий звонок) — состояние восстанавливается;
	\item сворачивание приложения — состояние сохраняется;
	\item слабое интернет-соединение — приложение работает автономно (данные локальные);
	\item низкий заряд батареи — приложение не потребляет избыточную энергию.
\end{itemize}

\subsection{Совместимость с версиями Android}

Приложение было протестировано на следующих версиях Android:
\begin{itemize}
	\item Android 7.0 (API 24) — эмулятор, работает корректно;
	\item Android 8.0 (API 26) — эмулятор, работает корректно;
	\item Android 9.0 (API 28) — физическое устройство, работает корректно;
	\item Android 10 (API 29) — физическое устройство, работает корректно;
	\item Android 11 (API 30) — физическое устройство, работает корректно;
	\item Android 12 (API 31) — физическое устройство, работает корректно;
	\item Android 13 (API 33) — эмулятор, работает корректно.
\end{itemize}

\subsection{Выявленные дефекты и их устранение}

В процессе тестирования были выявлены следующие дефекты:
\begin{enumerate}
	\item \textbf{Дефект 1}: При быстром переключении между вкладками приложение иногда зависало. \textbf{Устранение}: оптимизирована загрузка фрагментов, добавлена проверка на повторное создание.
	\item \textbf{Дефект 2}: Изображение аватара не загружалось при первом запуске. \textbf{Устранение}: добавлена загрузка изображения в фоновом потоке.
	\item \textbf{Дефект 3}: Уведомления не отправлялись на Android 12 и выше. \textbf{Устранение}: добавлена поддержка новых разрешений для уведомлений.
\end{enumerate}

\subsection{Выводы по тестированию}

В результате тестирования подтверждено, что разработанное приложение соответствует всем функциональным и нефункциональным требованиям. Приложение работает стабильно на всех тестовых устройствах и версиях Android. Выявленные дефекты были устранены. Приложение готово к эксплуатации.